\section*{Chapter \ref{ch:g11:absorbance} --- Colour and Absorbance}

\begin{solution}{exo:g11:absorbance:1}
Absorbing orange, it looks blue, the colour facing orange on the wheel.
Absorbing violet, it looks yellow.
\end{solution}

\begin{solution}{exo:g11:absorbance:2}
A green solution absorbs mainly red light, the colour facing green on
the wheel: from about \qty{620}{nm} to \qty{750}{nm}.
\end{solution}

\begin{solution}{exo:g11:absorbance:3}
Blue dye: $\lambda_{\max} = \qty{629}{nm}$, $A = 0.82$. Yellow dye:
$\lambda_{\max} = \qty{427}{nm}$, $A = 0.795$, about 0.80.
\end{solution}

\begin{solution}{exo:g11:absorbance:4}
$A$ is proportional to the concentration: three times more gives
$A = 0.90$; half as much gives $A = 0.15$.
\end{solution}

\begin{solution}{exo:g11:absorbance:5}
The blank is a cell filled with the pure \omterm{def:g6:solutions-and-solubility:solute}{solvent}. Setting $A = 0$ with it
removes what the cell, the \omterm{def:g6:solutions-and-solubility:solute}{solvent} and the instrument absorb, so that
the readings that follow are due to the \omterm{def:g6:solutions-and-solubility:solute}{solute} alone.
\end{solution}

\begin{solution}{exo:g11:absorbance:6}
$C = 0.41 / 0.164 = \qty{2.5}{mg/L}$.
\end{solution}

\begin{solution}{exo:g11:absorbance:7}
At \qty{629}{nm} the \omterm{def:g11:absorbance:absorbance}{absorbance} is largest, so the measurement is most
sensitive (small concentrations still give readable \omterm{def:g11:absorbance:absorbance}{absorbances}); and at
the top of the band the curve is flat, so a small error on the wavelength
hardly changes the reading. At \qty{550}{nm} the \omterm{def:g11:absorbance:absorbance}{absorbance} is small and
changes fast with the wavelength.
\end{solution}

\begin{solution}{exo:g11:absorbance:8}
Diluted: $C = 0.48 / 0.164 = \qty{2.9}{mg/L}$; the drink:
$5 \times 2.9 = \qty{15}{mg/L}$ (more precisely $14.6$).
\end{solution}

\begin{solution}{exo:g11:absorbance:9}
$a = A / (\ell\, C_m) = 0.795 / (1.00 \times 0.0150) = \qty{53.0}{L/(g.cm)}$;
$\varepsilon = a \times M = 53.0 \times 534.4 = \qty{2.83e4}{L/(mol.cm)}$.
\end{solution}

\begin{solution}{exo:g11:absorbance:10}
Above about \qty{520}{nm}. At \qty{629}{nm} the yellow dye does not
absorb at all, so the whole \omterm{def:g11:absorbance:absorbance}{absorbance} there comes from the blue dye.
\end{solution}

\begin{solution}{exo:g11:absorbance:11}
It doubles: $A = \varepsilon \ell c$ is proportional to $\ell$; the light
crosses twice as much solution, hence twice as many absorbing \omterm{def:g7:atoms-and-molecules:molecule}{molecules}.
\end{solution}

\begin{solution}{exo:g11:absorbance:12}
$A/C$: $0.164$, $0.162$, $0.115$, $0.0675$. The ratio is constant only
for the first two: beyond about \qty{10}{mg/L} (\omterm{def:g11:absorbance:absorbance}{absorbances} above about
1.6) the law no longer holds. Only the standards at 5 and \qty{10}{mg/L}
(and lower ones) may be used; an unknown reading 2.5 must be diluted, for
example five times, and measured again.
\end{solution}

\begin{solution}{exo:g11:absorbance:13}
Blue dye, from \qty{629}{nm} alone: $C = 0.574 / 0.164 = \qty{3.5}{mg/L}$.
Yellow dye, from \qty{427}{nm}: $C = 0.53 / 0.0530 = \qty{10}{mg/L}$.
\end{solution}

\begin{solution}{exo:g11:absorbance:14}
She finds $0.368 / 0.164 = \qty{2.24}{mg/L}$. The true \omterm{def:g11:absorbance:absorbance}{absorbance} is
$0.368 - 0.040 = 0.328$, so the true concentration is
$0.328 / 0.164 = \qty{2.00}{mg/L}$. She is \qty{12}{\percent} too high.
\end{solution}

\begin{solution}{exo:g11:absorbance:15}
$10 \times 60 = \qty{600}{mg}$ per day, in $600 / 20 = \qty{30}{L}$ of
drink. No one drinks thirty litres a day: the drink alone is no realistic
risk, although the dye may also come from other foods.
\end{solution}

\begin{solution}{pb:g11:absorbance:1}
\textbf{1.} Orange-red, the colour facing blue on the wheel.

\textbf{2.} $\lambda_{\max} = \qty{629}{nm}$.

\textbf{3.} At \qty{450}{nm} the light is blue: the dye lets it through,
which is exactly why the drink looks blue.

\textbf{4.} At \qty{629}{nm}.

\textbf{5.} $\qty{50.0}{mg} / \qty{0.5000}{L} = \qty{100}{mg/L}$.

\textbf{6.} \omterm{def:g10:concentration-and-dilution:dilution}{Dilution factor} $100 / 3.0 = 33.3$, so
$V_0 = 100.0 / 33.3 = \qty{3.00}{mL}$: a graduated pipette (or a
burette) and a \qty{100.0}{mL} volumetric flask.

\textbf{7.} $A/C$ = 0.168, 0.163, 0.165, 0.163, 0.165: all close to
0.164, so $A \approx 0.164\, C$, a line through the origin.

\textbf{8.} A solution without dye ($C = 0$) is the blank, set to
$A = 0$.

\textbf{9.} $a = \qty{0.164}{L/mg}$ per centimetre, that is
\qty{164}{L/(g.cm)}; $\varepsilon = 164 \times 792.9 =
\qty{1.30e5}{L/(mol.cm)}$.

\textbf{10.} $F = 50.0 / 25.0 = 2$.

\textbf{11.} $C = 0.590 / 0.164 = \qty{3.60}{mg/L}$.

\textbf{12.} $2 \times 3.60 = \qty{7.20}{mg/L}$.

\textbf{13.} About $2 \times 0.590 = 1.18$, above the highest standard
(0.823): the reading would lie outside the calibrated range, where the
law may fail. Diluting brings it among the standards.

\textbf{14.} $\qty{7.20}{mg/L} \times \qty{0.500}{L} = \qty{3.60}{mg}$.

\textbf{15.} $n = \num{3.60e-3} / 792.9 = \qty{4.54e-6}{mol}$.

\textbf{16.} $6 \times 30 = \qty{180}{mg}$ per day.

\textbf{17.} $180 / 3.60 = 50$ bottles.

\textbf{18.} $50 \times 0.500 = \qty{25}{L}$ in a day: impossible to
drink. The dye of this drink alone cannot bring a child near the limit,
though dyes from several foods add up.

\textbf{19.} $6 \times 70 = \qty{420}{mg}$, that is $420 / 3.60 \approx
117$ bottles.

\textbf{20.} About 50 bottles of \qty{500}{mL}, \qty{25}{L} of drink, in
a single day.
\end{solution}
